\subsection{Continuity with the standard formalism and expected corrections}
  \label{subsec:continuity-with-standard-formalism}

  The emergence of the Schr\"{o}dinger equation in the projected weak-amplitude regime
  explains why standard quantum mechanics remains empirically valid across all currently
  accessible laboratory domains.
  In those regimes, the effective description is sufficiently far from the saturation threshold
  that the nonlinear corrections induced by bounded projective capacity remain negligible.

  The present framework nevertheless predicts that exact linearity cannot hold arbitrarily.
  When localization becomes too sharp, when relational gradients become too large, or when
  the projective description approaches a saturation boundary, the effective Schr\"{o}dinger
  dynamics must receive corrections.
  These corrections preserve the ordinary formalism in the infrared and moderate-resolution
  regime, while producing departures only near the structural limits of projectability.

  Cosmochrony does not reinterpret the Schr\"{o}dinger equation as false or merely
  approximate in an informal sense.
  Rather, it identifies the precise structural reason why it must arise, and also why it cannot
  be ultimate.
  The quantum formalism is the natural low-resolution description of a deeper relational
  dynamics, just as effective geometry is the low-curvature description of a deeper projective
  structure.

  In this sense, the relation between Cosmochrony and quantum mechanics is analogous to
  the relation between kinetic theory and hydrodynamics.
  The effective equation is real, predictive, and operationally complete within its domain.
  Its domain of validity is, however, determined by a more primitive substrate.
  The wavefunction $\psi$ is therefore not eliminated, but reinterpreted:
  it remains the correct effective descriptor of projected modal structure within the
  projected-stable sector, while its apparent fundamentality is replaced by an emergent
  status grounded in the joint structural constraints of bounded substrate capacity and
  finite projective resolution.
